\documentclass[12pt,epsbox]{article}
\usepackage{authblk}
\usepackage{hyperref}
\usepackage{color}
\usepackage[dvips]{graphicx}
\usepackage{amsmath}


\textwidth=170true mm \textheight=220true mm
\oddsidemargin=-5true mm\evensidemargin=-5true mm
\topmargin=-15true mm
\linespread{1.1}

\title{Notes on the Exact RG equation and the Wheeler-DeWitt equation}
%%%%%%%%%%%%%%%%%%%%%%%%%%%%%%%%%%%%%%%%%%%%%%%%%%%%%%%%%%%%%%%%%%%%
%% Math symbols etc
\newcommand{\rom}[1]{\mathrm{#1}}

% latex equations

\newcommand{\be}{\begin{equation}}
\newcommand{\ee}{\end{equation}}
\newcommand{\bs}{\begin{split}}
\newcommand{\es}{\end{split}}
\newcommand{\bea}{\begin{eqnarray}}
\newcommand{\eea}{\end{eqnarray}}
\newcommand{\ba}{\begin{eqnarray}}
\newcommand{\ea}{\end{eqnarray}}
\newcommand{\nn}{\nonumber \\}
\newcommand{\eqn}[1]{(\ref{#1})}
\newcommand{\beq}{\begin{equation}}
\newcommand{\eeq}{\end{equation}}
\newcommand{\beqa}{\begin{eqnarray}}
\newcommand{\eeqa}{\end{eqnarray}}
\newcommand{\beqar}{\begin{eqnarray*}}
\newcommand{\eeqar}{\end{eqnarray*}}
\numberwithin{equation}{section}
%\renewcommand{\theequation}{\thesection.\arabic{equation}}

%%%%%%%%%%%%%%%%%%%%%%%%%%%%%%%%%%%%%%%%%%%%%%%%%%%%%%%%%%%%%%%%%%%%
% useful stuff

\newcommand{\labell}[1]{\label{#1}} %\qquad_{#1}} %{\label{#1}}
\newcommand{\labels}[1]{\label{#1}} %{\vskip-2ex$_{#1}$\label{#1}}
\newcommand{\reef}[1]{(\ref{#1})}
\newcommand{\ssc}{\scriptscriptstyle}
\newcommand{\eg}{{\it e.g.,}\ }
\newcommand{\ie}{{\it i.e.,}\ }
\newcommand{\n}{\nu}
\newcommand{\m}{\mu}
\newcommand{\mo}{\mathcal{O}}
\newcommand{\e}{\epsilon}

%%%%%%%%%%%%%%%%%%%%%%%%%%%%%%%%%%%%%%%%%%%%%%%%%%%%%%%%%%%%%%%%%%%%

\newcommand{\gsim}{\mathrel{\raisebox{-.6ex}{$\stackrel{\textstyle\rangle}{\si
m}$}}}
\newcommand{\lsim}{\mathrel{\raisebox{-.6ex}{$\stackrel{\textstyle<}{\si
m}$}}}

%%%%%%%%%%%%%%%%%%%%%%%%%%%%%%%%%%%%%%%%%%%%%%%%%%%%%%%%%%%%%%%%%%%%

\begin{document}

  \begin{flushleft}
    {\bf OUJ-FTC-11}
  \end{flushleft}
  
  \begin{center}
    \vspace{37pt}
    {\LARGE Notes on the Exact RG equation and the Wheeler-DeWitt equation}\\[1em]
    \vspace{17pt}
    {\large Hideto Kamei\footnote{hidetokamei@gmail.com}}\\[0.5em]
    \vspace{17pt}
    {Field theory group, \\
  The Open University of Japan,\\
  Chiba 261-8586, Japan}\\[1em]
    \vspace{17pt}
    {\large February 25, 2024}
  \end{center}
  \vspace{17pt}

In this note, in the context of the AdS/CFT correspondence, the holographic derivation of the Wilsonian effective action is proposed. Then, the exact RG equation in the boundary theory is derived from the Wheeler-DeWitt equation of the bulk, following the suggestion of \cite{dVV},\cite{VV}, and\cite{Polchinski}. The relationship between the exact RG and Stochastic Quantization\cite{Parisi},\cite{Damgaard} is briefly discussed.

\section{Introduction}
Here in this note, we would like to promote Holographic RG to exact RG, which corresponds to quantum gravity in the bulk.
By Verlinde et al \cite{dVV}, it was shown that the RG flow equation of the CFT can be derived from classical Equation of motion of the bulk\footnote{For a review, please refer to \cite{Sakai}.}. We generalize the RG flow equation to exact RG, which is derived from Wheeler-Dewitt equation in the bulk. We identify the expression of the boundary action and counterterm is properly subtracted in the derivation.

Now we outline the structure of this note. In section 2, the setup of the derivation is given. In section 3, the role of the boundary condition in the setup is explained. In section 4, procedure for splitting the bulk action and identifying the boundary action is given. In section 5, derivation of the exact RG equation from the bulk is given. In section 6, interpretation of the equation as the exact RG is given. In section 7, an example to show the validity of the methodology is given. In section 8, discussion on the result and possible future research is given. In section 9, update on research since this note was first written is given.

%\section{Review of Holographic RG}

%we will leave the reader to refer to \cite{dVV}, or \cite{Sakai}, or my master thesis. %\cite{HKamei}
%Necessary material to understand this note will be,
%\begin{itemize}
%\item About AdS/CFT correspondence
%\item Hamilton-Jacobi Equation
%\item extrinsic curvature and decomposition of the Ricci scalar
%\item deriving RG flow equation of the CFT side from classical EOM in the AdS side
%\end{itemize}

\section{Deriving the exact RG from the bulk theory }

%We could derive Wilsonian RG equation from the bulk equation.
%For that purpose, we wish to solve Wheeler-de Witt equation of the bulk, mostly following \cite{dVV}. we suspect it will give us Wilsonian renormalization equation, in the form of \cite{Morris}, (3.16).

Here we wish to derive the exact renormalization group equation (In the form of \cite{Morris},\cite{Morris2},\cite{Morris3},\cite{Oliver}) of $\mathcal{N}=4$ SYM from the Wheeler-DeWitt equation in the AdS. Most of the discussion follows from the beautiful paper on the holographic renormalization, \cite{dVV}. In the paper, the Callan-Symanzik type of the RG flow equation was derived from the condition that the Hamiltonian is zero, $H=0$, and the Wheeler-DeWitt equation is the quantum extension of $H=0$.

We assume that the bulk action contains gravity and a scalar field (the extension to including many scalar fields and including other fields is straightforward).
%\begin{equation}
%S= \int \left( {}^5 R + V(\phi) + \frac{1}{2} \partial_{\mu} \phi \partial^{\mu} \phi \right) \sqrt{g} d^5x.
%\end{equation}
We take the parametrization of the 5d metric as,
\begin{equation}
ds^2 = g_{\mu\nu} dx^{\mu}dx^{\nu} = N^2 dr^2 +h_{ij}(dx^i +N^i dr)(dx^j +N^j dr).
\end{equation}

We will work in the Euclidean signature everywhere in this note.
Here we fix the gauge of the metric as $N=1, N_i=0$. %(The coefficient of the metric part and the scalar field part will look similar and look better if we take $N=2$ though) <- not sure
 We decompose the Ricci scalar using extrinsic curvature, %as explained in \cite{Wald} <- textbook.
\begin{equation}
K_{ij} \equiv \frac{1}{2N} \left( \partial_r h_{ij} -D_i N_j -D_j N_i \right)=\frac{1}{2}\left( \partial_r h_{ij} \right).
\end{equation}

From now on we treat the metric as fields and use the condensed notation, $h_{ij} \equiv h^I ,K_{ij} \equiv K^I$.
Further we will use so-called DeWitt's supermetric $G_{IJ} \equiv G^{ijkl} \equiv \sqrt{h}\left( \frac{1}{2} (h^{ik}h^{jl}+h^{il}h^{jk})-h^{ij}h^{kl} \right)$ and its inverse $G^{IJ}$ for the condensed notation of the action.
We can define the conjugate momentum of the metric and the scalar field as $\pi_{I}=\delta \mathcal{L} / \delta \partial_r h^I$,
and $p=\delta \mathcal{L} / \delta \partial_r \phi$, and construct the Hamiltonian $H=\pi_I \partial_r h^I + p \partial_r \phi-\mathcal{L}$.
%Hamiltonian of the theory is,
%\begin{equation}
%\begin{split}
%H &= \frac{1}{\sqrt{h}} (h_{ki}h_{jl}-\frac{1}{3}h_{ij}h_{kl})\pi^{ij} \pi^{kl}+ \frac{1}{2} \frac{1}{\sqrt{h}} p^2 -\sqrt{h}\left( {}^4R-V(\phi)
%- \frac{1}{2} (\partial_i \phi)^2 \right).\\
%  &= G^{IJ} \pi_I \pi_J + \frac{1}{2\sqrt{h}}  p^2 - \sqrt{h}\left( {}^4R-V(\phi)
%- \frac{1}{2} (\partial_i \phi)^2 \right).
%\end{split}
%\end{equation}
%Classically, $H=0$ (which comes from Equation of Motion for radial component of the metric, $N$), but here we wish to quantize the theory.
The equation we wish to solve is the Wheeler-DeWitt equation, which is $H \Psi =0$. Analogous to the Shr\"odinger equation, we simply have to replace momentum $\pi_I$, $p$ with the derivative $\delta / \delta h^I$, $\delta / \delta \phi$. We wish to proceed following \cite{VV} to derive the exact RG equation of the CFT holgraphically.
%We could proceed and derive the exact RG for changing the cutoff of the theory, following \cite{dVV}, but here we don't explain it in detail.
In the classical case, we will get the expression of the radial velocity of the fields (with respect to the 4d action) from Hamilton-Jacobi equation,

%We can think of $\Psi$ as the partition function of the bulk theory inside $r=const$ hypersurface that we are considering, and at the same time partition function of the CFT  below some cutoff scale (IR). We can therefore write $\Psi=\exp(-S)$, where $S$ is the effective action of the CFT plus divergent part, which is always subtracted by counterterm when we consider field theory. We write down this divergent part again \cite{dVV}. They perform the derivative counting, and write down the counterterm as having up to two derivatives of the boundary theory.
%It can be understood, from $h_{ij} \sim \exp(2r)$, and we need to use the inverse of the metric to contract the indices for each derivative, therefore making the quantity much smaller.

%Therefore the most general divergent term is,

%\begin{equation}
%S_{div}= S^{(0)}_{div}+S^{(2)}_{div} = \int \sqrt{h} U(\phi) d^4x + \int \sqrt{h} (\Phi (\phi)R + \frac{1}{2} L(\phi) \partial_i \phi \partial^i \phi) d^4x.
%\end{equation}
%If we plug in $\Psi= \exp(-S)$ ($S= S_{div}+\Gamma$, where $\Gamma$ is the effective action of the CFT) to Wheeler-DeWitt equation, we will have
%\begin{equation}
%\left( G^{IJ} \frac{\delta }{\delta h^I}\frac{\delta }{\delta h^J} + \frac{\delta^2 }{\delta \phi^2}
%-\sqrt{h}({}^4R-V-\partial_i \phi \partial^i \phi) \right)\Psi=0.
%\end{equation}

%Then, we wish to solve it by separating the equation into different number of derivatives.
%The part up to two derivatives will give
%\begin{equation}
%G^{IJ}\left( \frac{\delta S_{div}}{\delta h^I} \frac{\delta S_{div}}{\delta h^J} - \frac{\delta^2 S_{div}}{\delta h^I \delta h^J} \right)
%+ \frac{1}{2} \left( \frac{\delta S_{div}}{\delta \phi} \frac{\delta S_{div}}{\delta \phi} - \frac{\delta^2 S_{div}}{\delta \phi^2} \right)
%= \sqrt{h} (R-V- \frac{1}{2}(\partial \phi)^2 ).
%\end{equation}
%Note here that $(\delta S_{div}/ \delta \phi)^2 \gg \delta^2 S_{div} / \delta \phi^2 $, simply because $S_{div}$ is big. We could also understand this as coming from $1/N$ effect, since Newton constant in the action includes $N$.
%>> latex de doyatte kaku?
%This form of the 5d potential strongly suggest that we can construct this AdS/CFT setting from stochastic quantization, starting from
%$\partial_r \phi= (1/2) (\delta S_{div} / \delta \phi) + \eta$, and $K^I = G^{IJ} (\delta S_{div} / \delta h^J) + \eta^a E^I_a$,
% RHS no fugo ga chigau you ni mieruga, r ga sotomuki nanode tabun OK.
% where $E$ is the inverse of the vielbein from DeWitt's supermetric. (See \cite{Susanne} for the construction of the gravitational theory out of Stochastic quantization)

%We wish to identify the terms with more than 4 derivatives as the Exact RG equation.
%\begin{equation}
%\begin{split}
%G^{IJ}\left( 2\frac{\delta S_{div}}{\delta h^I}\frac{\delta \Gamma}{\delta h^J}- \frac{\delta^2 \Gamma}{\delta h^I \delta h^J}
%+ \frac{\delta \Gamma}{\delta h^I} \frac{\delta \Gamma}{\delta h^J} \right)+ \frac{1}{2}\left( 2\frac{\delta S_{div}}{\delta \phi} \frac{\delta \Gamma}{\delta \phi}- \frac{\delta^2 \Gamma}{\delta \phi \delta \phi}
%+ \frac{\delta \Gamma}{\delta \phi} \frac{\delta \Gamma}{\delta \phi} \right)\\
%= -G^{IJ}\frac{\delta S^{(2)}_{div}}{\delta h^I}\frac{\delta S^{(2)}_{div}}{\delta h^J}-\frac{1}{2}\left( \frac{\delta S^{(2)}_{div}}{\delta \phi} \right)^2.
%\end{split}
%\end{equation}

%Together with that equation,
%We wish to identify the radial velocity of the fields as Following \cite{VV},
\begin{equation}
\begin{split}
\frac{\delta S_{UV}}{\delta \phi} &= \partial_r \phi\\
\frac{\delta S_{UV}}{\delta h^I}  &=  \frac{1}{16\pi G_N} G_{IJ} \partial_r h^{J} \label{velocity},
\end{split}
\end{equation}
and in the quantum case the same equation holds for field redefinition, but we will need a correction term $\mathcal{F}(\phi)$ as in (\ref{fieldtransf}), which will be consistent with Wheeler-DeWitt equation (\ref{HERG1}) we will derive.
%collect the terms more than 4 derivatives to get the Exact RG with respect to the cutoff.


% (It is basically Hamilton-Jacobi equation for the divergent part),
%\begin{equation}
%\frac{d}{dr} \Gamma = 2G^{IJ} \frac{\delta S_{div}}{\delta h^I} \frac{\delta \Gamma}{\delta h^J}
%+ \frac{\delta S_{div}}{\delta \phi}\frac{\delta \Gamma}{\delta \phi}
%\end{equation}
%and if we plug it into part of Wheeler-DeWitt with more than 4 derivatives, we will get
%\begin{equation}
%-\frac{d}{dr} \Gamma = G^{IJ}\left( \frac{\delta \Gamma}{\delta h^I} \frac{\delta \Gamma}{\delta h^J} -\frac{\delta^2 \Gamma}{\delta h^I \delta h^J}  \right)
%+ \frac{1}{2}\left( \frac{\delta \Gamma}{\delta \phi} \frac{\delta \Gamma}{\delta \phi} -\frac{\delta^2 \Gamma}{\delta \phi \delta \phi} \right) +G^{IJ}\frac{\delta S^{(2)}_{div}}{\delta h^I}\frac{\delta S^{(2)}_{div}}{\delta h^J} + \frac{1}{2}\left( \frac{\delta S^{(2)}_{div}}{\delta \phi} \right)^2.
%\end{equation}

%This is the form of the RG equation of Polchinski type. If we recover 5d Newton constant $G$ and compare terms in RHS, the first term will have power of $G$ and much smaller than the rest of the terms.
 %Moreover, if we take the flat space limit, the last term, which is responsible for anomaly, also drops out.

%Let's discuss this equation as the exact RG. We can actually see that it is an exact RG by, first considering "course graining function"\cite{Oliver}. We can consider, when we change the physical cutoff scale from $t_0$ ($t$ is defined by $dt \equiv -d\Lambda/ \Lambda$) to $t=t_0+dt$, the course graining function as $b_{dt}(\phi_0)=\phi_0- (1/2) \left( \delta S_0 / \delta \phi_0 \right) dt$. The effective action after "course graining" will be then,
%\begin{equation}
%e^{-S}= \int \mathcal{D}\phi_0 \  \delta (\phi- b_{dt}(\phi_0)) \ e^{-S_0(\phi_0)}
%\end{equation}
%Then, we will have
%\begin{equation}
%\mathcal{Z}= \int \mathcal{D}\phi \ e^{-S} =\int \mathcal{D}\phi_0 \ e^{-S_0(\phi_0)}
%\end{equation}
%as it should be. Using then $\phi=\phi_0-  (1/2)( \delta S_0 / \delta \phi_0 ) dt $, and expanding $\mathcal{D}\phi$ and $\exp(-S)$ up to first order of $dt$, we will get

%\begin{equation}
%\begin{split}
%\mathcal{Z} &= \int \mathcal{D} \phi \ e^{-S}\\
%            &\sim \int \mathcal{D}\phi_0 \ \exp \left(-S_0+ \left( \left( \frac{\delta S_0}{\delta \phi_0} \right)^2 -\frac{\delta^2 S_0}{\delta \phi_0^2} \right) dt \right)
%\end{split}
%\end{equation}
%Which is precisely the final equation of the Holographic RG.
%The above "course graining" function is reasonable, since it will approach classical configuration $(\delta S / \delta \phi)=0$ for large scale. we suspect that the "course graining function" should satisfy detailed balance condition, when translated into the language of Fokker-Planck equation.


%It is the form of the generalized Holographic RG equation which was proposed in \cite{Miao}, obtained by the boundary consideration. It was shown to be equivalent to "weak form of Polchinski equation" in which the average of the path integral is taken. (Comparing Polchinski equation in $\mathcal{N}=4$, with this quantum version of the holographic equation, will be certainly a work which is remaining.)

%It is, if we define the partition function of this $\Gamma$ as $Z= \exp(-\Gamma)$, and define the momentum cutoff scale as $\Lambda=e^r$,
%\begin{equation}
%\Lambda  \frac{\partial} {\partial \Lambda}Z= \left( G^{IJ} \frac{\delta}{\delta h^I} \frac{\delta }{\delta h^J}
%+ \frac{1}{2} \frac{\delta}{\delta \phi^2}  \right) Z
%\end{equation}
%which is the Wilsonian RG equation in the form of \cite{Morris}
%(However, the $\Lambda$ factor in front of this equation is not there in \cite{Morris}, we need to find out why. also the sign on the RHS is also problem?)

\section{The role of the boundary condition and the fluctuation}

Here in this section, we would like to explain the role of the boundary condition in terms of the fluctuation in the boundary CFT, in order to motivate the holographic definition of the Wilsonian effective action in the following section.

First of all, we would like to claim that the boundary value of the bulk fields themselves could be considered as the dynamical fields in the boundary CFT in some cases. The main example we consider in this note is a scalar field in the AdS within the mass range $-d^2/4 < m^2 < -(d^2/4)+1$.
In this mass range, it is known, in addition to the standard Dirichlet boundary condition, we can also impose free boundary condition\cite{Witten2},\cite{Muck},\cite{Witten},\cite{WittenNeumann},\cite{MukundHRG}, and it is called "alternative quantization". Those two different boundary conditions correspond to different CFTs on the boundary, $CFT_D$ and $CFT_F$ respectively. In the case of free boundary condition, we also need to integrate over the value of the boundary condition of the bulk field, and therefore, the partition function of the $CFT_F$ can be obtained by promoting the source in the $CFT_D$ to a dynamical field and integrating over it as
\begin{equation}
\mathcal{Z}_F =\int \mathcal{D}J \mathcal{Z}_D = \int \mathcal{D}J \exp(-\Gamma[J])\label{FD}.
\end{equation}

This expression can be interpreted as the path integral formally (or statistical partition function if one wishes), where $J$ is now a dynamical field, and weight is $\exp(-\Gamma[J])$, $\Gamma$ being the generating functional for $CFT_D$. Therefore, in the case of the free boundary condition, the boundary value of the bulk field can be considered as the field in the boundary CFT itself, and we can consider an exact RG with respect to that.\footnote{Necessity of the free boundary condition was first realized by the necessity of having operator with small conformal dimension, which can't be realized by the naive Dirichlet boundary condition. Later it was realized those two CFTs are connected by an RG flow, by adding relevant perturbation $f \mathcal{O}^2$, where $\mathcal{O}$ is an operator to couple the source. (\ref{FD}) was obtained using that relationship.} In addition, also for other fields, it is known that in the free boundary condition (or equivalently Neumann boundary condition in the classical limit), the bulk fields on the boundary are interpreted as the dynamical fields in the boundary CFT.\cite{MukundOfer},\cite{Ross}.

%Note that there are two asymptotics of the scalar field at the boundary, $\phi \sim z^{\Delta_+} \alpha(x) + z^{\Delta_-} \beta(x)$, $\Delta_{\pm}$ being two solutions of $\Delta(\Delta-d)=m^2$ with $\Delta_+ > \Delta_-$. It is known, in the mass range, two different boundary conditions can be imposed. One is identifying faster-falling off mode at the boundary $\alpha(x)$ as the source in the CFT $J_{\alpha}$, and the other is identifying slower-falling off mode $\beta(x)$ as the source $J_{\beta}$. We call them as $\alpha$ theory and $\beta$ theory respectively, following \cite{WittenNeumann}. Those two boundary conditions corresponds to two different CFTs on the boundary. There is a relationship between the partition functions for the two CFTs, due to the fact that they are connected by an RG flow.\footnote{There is an RG flow from $\alpha$ theory to $\beta$ theory. We add a perturbation $\int f \mathcal{O}^2_{\alpha}$ to $\alpha$ theory, where $\mathcal{O}_{\alpha}$ is the field which couples to the source $J_{\alpha}$, and $f$ is a constant. IR fixed point can be realized by sending f to infinity, since it was a relevant perturbation, and there we will have $\beta$ theory.} The relationship for the partition functions is, namely, the partition function of the $\beta$ theory is equal to the partition function of the $\alpha$ theory where we integrate over $J_{\alpha}$. (That is, $\mathcal{Z}_{\beta}=\int \mathcal{D}J_{\alpha} \ \mathcal{Z}_{\alpha}[J_{\alpha}]$.)
%We could also interpret the relation as, the $\beta$ theory is equivalent to the $\alpha$ theory, where we promote Dirichlet boundary condition to the free boundary condition, in which we integrate over all the possible value of the boundary condition. That expression formally looks like path integral, where $J_{\alpha}$ is a dynamical field in the CFT and we integrate with the weight of $\mathcal{Z}_{\alpha}=\exp(-\Gamma[J_{\alpha}])$, $\Gamma$ being the generating functional. Therefore, in this case, free boundary condition\footnote{In the classical limit of the bulk, since the boundary value should satisfy $\delta S_{bulk} / \delta \phi_{bdry}=0$ from the equation of motion and using Hamilton-Jacobi equation to identify this as the momentum $p$, we will get the Neumann boundary condition, $\partial_r \phi=0$.} naturally arise. %even if we consider only Dirichlet boundary condition
%In terms of bulk path integral, the same expression would be obtained by first integrating over all the bulk fields except the boundary condition.

Another example for the free boundary condition is, trailing string in the AdS\cite{StochasticdeBoer},\cite{StochasticSon}. We consider a string stretching from the horizon of the black brane to the boundary at the finite radial distance. The endpoint on the boundary is identified as the position of a particle in the CFT. Since the boundary is at the finite radial distance, the fluctuation of the string in the bulk can be transferred to the fluctuation of the endpoint, which is the fluctuation of the particle in the CFT. Therefore, we need to sum over all the possible value of the boundary position of the string, and it is formally the same as the path integral for the position of the particle in the CFT. Therefore, in the case of the trailing string, we also take the free boundary condition, and the boundary condition is promoted to a dynamical field in the boundary. we expect, if we take the boundary at the finite radial distance\cite{VV}, boundary value of the closed string modes in the bulk might be also identified as the fields of the closed string modes of the boundary theory.

Therefore, we claimed that, in some cases in the AdS/CFT, free boundary condition naturally arise and the integration over the boundary condition could be formally identified as the path integral in the boundary CFT. The above discussion will not be precise enough, but hopefully enough to motivate the following holographic definition of the Wilsonian effective action.

%In this section, in order to motivate the following definition of the Wilsonian effective action of the CFT side from the bulk, we would like to claim that the bulk fields might be directly interpreted as the fields on the boundary as well. In usual AdS/CFT correspondence, Dirichlet boundary condition is taken and the value of the bulk field on the boundary is identified as the value of the source. However, when we take Neumann boundary condition, the value of the bulk fields on the boundary are considered to be dynamical field and it should be integrated over\cite{StochasticSon},\cite{MukundHRG}\footnote{We can take Neumann boundary condition for a scalar field with some particular range of mass in the bulk, $-d^2/4 < m^2 < -(d^2/4)+1$,\cite{Witten},\cite{Witten2},\cite{Muck}, and it is called "alternative quantization". We can also take the Neumann boundary condition when we put the boundary at some regularized hypersurface, as in the case of stochastic string\cite{StochasticSon},\cite{StochasticdeBoer}.}.
%the value of the bulk field is identified as the value of the one point function. Further, if we consider the fluctuation in AdS/CFT,
%it is known that we need to integrate over all the possible value of boundary condition
%In the case of the stochastic string, where we consider a fundamental string stretching from the regularized boundary to the horizon, and identify the boundary endpoint as the coordinate of a particle, the integration over the boundary condition could be thought of as the path integral of the boundary theory with respect to the coordinate of the particle. When we write the partition function in the bulk and perform the path integral for the bulk fields except for the fields on the boundary (remaining sum is over the boundary condition), the expression of the total partition function $\mathcal{Z}=\int \mathcal{D}\phi_{bdry} \exp (-S_{bulk}(\phi_{bdry}))$ generally looks like the path integral for the boundary theory, in which the value of the Neumann boundary condition would be interpreted as the value of the field which needs to be summed over. This discussion is far from complete, but hopefully enough to motivate the following definition of the Wilsonian effective action of the CFT side from the bulk. In the above discussion the boundary is at some large r which corresponds to the fundamental scale of the CFT, but we will move this boundary to derive the exact RG equation.

\section{The definition of the boundary Wilsonian effective action from the bulk}

When we define the Wilsonian effective action of the CFT side, we observe the CFT side with respect to a particular energy scale. In the context of AdS/CFT, it would be natural to concentrate on a hypersurface in the AdS, which will manifest the physics of the CFT side with respect to a particular energy scale we observe the physics \cite{dVV}. There the radial coordinate in the AdS is related to the energy scale of the observation in the CFT side. %The first step for defining the Wilsonian effective action of CFT side from the bulk comes from specifying the hypersurface in the AdS, and
Let's take the hypersurface as $r=const$. Then, following \cite{Polchinski}, the total partition function of the bulk, which is also equal to the total partition of the CFT, can be written as the path integral over the fields on this boundary, and fields in the region of larger $r$ than this hypersurface (I denote it as UV, since it corresponds to the UV degree of freedom in the CFT side), and fields in the smaller $r$ region (I denote it as IR). Then, we can separate the expression of the bulk gravity action into UV part and IR part, since there is no cross term in the action for the fields with different $r$. Note that they are still the function of the Dirichlet boundary condition on the hypersurface.

 \begin{equation}
 \begin{split}
 \mathcal{Z} &= \int \mathcal{D}\phi_{UV} \mathcal{D} \phi_{IR} \mathcal{D} \phi_r \exp(-S^{bulk}_{UV}) \exp(-S^{bulk}_{IR})\\
             &= \int \mathcal{D}\phi_r \Psi_{UV}(r,\phi_r) \Psi_{IR}(r,\phi_r)\\
             &= \int \mathcal{D}\phi_r \exp(-S(r,\phi_r))\label{WilsonianPartition}
 \end{split}
 \end{equation}
 We can first perform the path integral for the fields of the UV and the IR part, and it will give the result of the second line. We expressed the UV part of the partition function, or wave function as $\Psi_{UV}(r,\phi_r)= \exp(-S^{bulk}_{UV}(r,\phi))$, and IR part of the wave function as $\Psi_{IR}(r,\phi_r)= \exp(-S^{bulk}_{IR}(r,\phi))$. Further, we can define the total low energy effective action as\\
 $\exp(-S(r,\phi_r)) \equiv \Psi_{UV}(r,\phi_r)\Psi_{IR}(r,\phi_r)$, which is the function of the boundary condition at $r$. The last line of (\ref{WilsonianPartition}) looks like the expression of the CFT partition function, with respect to the effective field $\phi_r$ and Wilsonian effective action $S(r,\phi_r)$, and it is tempting to define the Wilsonian effective action at some energy scale, which is determined by $r$, as simply this $S(r,\phi_r)$, holographically. Indeed, defining this way, we can show this "Wilsonian effective action" will follow the Exact RG equation with respect to the change of the energy scale, which is the boundary position $r$ in the bulk. This is given by solving the Wheeler-DeWitt equation in the bulk, which is the quantum extension of the condition that the Hamiltonian is zero for the gravity. Also note, that the above definition of the Wilsonian effective action, reduces to the definition of \cite{MukundHRG} in the bulk classical limit, as it should.

\section{Holographic Exact RG at finite radial distance}


%From these equations, we will have
%\begin{equation}
%\begin{split}
%& G^{IJ}\left( \frac{\delta \Gamma}{\delta h^I} \frac{\delta \Gamma}{\delta h^J} -\frac{\delta^2 \Gamma}{\delta h^I \delta h^J}  \right)
%+ \frac{1}{2}\left( \frac{\delta \Gamma}{\delta \phi} \frac{\delta \Gamma}{\delta \phi} -\frac{\delta^2 \Gamma}{\delta \phi \delta \phi} \right)\\
%=& \sqrt{h}({}^4R-V-\partial_i \phi \partial^i \phi)\\
%=& G^{IJ}\left( \frac{\delta S_{div}}{\delta h^I} \frac{\delta S_{div}}{\delta h^J} -\frac{\delta^2 S_{div}}{\delta h^I \delta h^J}  \right)
%+ \frac{1}{2}\left( \frac{\delta S_{div}}{\delta \phi} \frac{\delta S_{div}}{\delta \phi} -\frac{\delta^2 S_{div}}{\delta \phi \delta \phi} \right)
%\end{split}
%\end{equation}
%from these equation, together with (\ref{velocity}), and parametrization $dr= d\Lambda / \Lambda$ we will get
%\begin{equation}
%\Lambda \partial_{\Lambda} S = G^{IJ}\left( \frac{\delta S}{\delta h^I} \frac{\delta S}{\delta h^J} -\frac{\delta^2 S}{\delta h^I \delta h^J}  \right)
%+ \frac{1}{2}\left( \frac{\delta S}{\delta \phi} \frac{\delta S}{\delta \phi} -\frac{\delta^2 S}{\delta \phi \delta \phi} \right)
%\end{equation}
%which is exact RG equation for S.

Let's then derive the exact RG equation from the WDW equation in the bulk following the spirit of \cite{VV}. To make a contact with Polchinski type\cite{PolchinskiRG} of the Exact RG, let's first consider the regime where we can treat the metric as the classical background, and concentrate on the quantized scalar field in the AdS. This is done by recovering the Newton's constant $G_N$ to the Wheeler-DeWitt equation, and expanding it with $G_N$. This is done following \cite{GiladSchwinger}.

The Wheeler-DeWitt equation is,
\begin{equation}
\left[ 16 \pi G_N G^{IJ} \frac{\delta^2}{\delta h^I \delta h^J} - \frac{\sqrt{h}}{16\pi G_N} \left( (R-V) + (16\pi G_N) H_{matter} \right)  \right] \Psi=0.
\end{equation}

Both UV and IR wavefunction, $\Psi_{UV}$ and $\Psi_{IR}$ are expected to satisfy the above WDW equation separately.
Here we assume that $V$ is a constant (which doesn't depend on $\phi$), and $H_{matter}= (p^2/2) + W(\phi) -\Phi(\phi)R + (1/2)(\partial_i \phi)^2$, where $W(\phi)$ is the $\phi$ dependent part of the 5d potential, and its order $\phi^2$ has the coefficient $m^2/2$ which satisfies the bound that we discussed. We assume that $H_{matter}$ is the order of $\mathcal{O}(G_N^0)$. We further assume that we can expand
$S=-\log \Psi$ as $S=(S^0 / 16 \pi G_N) + S^1 + \mathcal{O}(G_N^1)$, and suppose $S_0$ depends only on the metric, and not on the scalar field. By expanding the Wheeler-DeWitt equation for $\mathcal{O}(G_N^{-1})$, while still seperating $S_0$ into the UV part and the IR part with respect to the cutoff, we get

\begin{equation}
\frac{1}{16 \pi G_N} \left(\frac{\delta S^0_{UV}}{\delta h} \right)^2 - \left( \frac{\delta^2 S^0_{UV}}{\delta h^2} \right) = -\frac{1}{16 \pi G_N} (\sqrt{h}) (V-R) =
\frac{1}{16 \pi G_N} \left(\frac{\delta S^0_{IR}}{\delta h} \right)^2 - \left( \frac{\delta^2 S^0_{IR}}{\delta h^2} \right).
\end{equation}

Note that the indices of the metric $h^I$ are implicitly contracted in the above expression. Strictly speaking, $\delta^2 S / \delta h^2$ term is order $\mathcal{O}(G_N^0)$, but we added it for simplicity of the next order \cite{GiladSchwinger}. This form of the 5d potential implies that the bulk gravity could be constructed out of the Stochastic Quantization\footnote{Please refer to \cite{Nishioka} for an analogous discussion.}, $ -\partial_r h^I=-G^{IJ}(\delta S_{UV} / \delta h^J) + E^I_a \xi^a$, where $\xi_a$ is the gaussian noise and $E$ is the inverse of the vielbein for $G_{IJ}$\cite{Susanne}. For $\mathcal{O}(G_N^0)$, we get
\begin{equation}
\begin{split}
G^{IJ} \frac{\delta S^0_{UV}}{\delta h^I} \frac{\delta S^1_{UV}}{\delta h^J} +\frac{1}{2} \left( \frac{ \delta S^1_{UV}}{\delta \phi} \right)^2 - \frac{1}{2}\frac{\delta^2 S^1_{UV}}{\delta \phi^2} &= -(\sqrt{h}) (W(\phi)-\Phi(\phi)R+\frac{1}{2}(\partial_i \phi)^2)\\
=&G^{IJ}\frac{\delta S^0_{IR}}{\delta h^I}\frac{\delta S^1_{IR}}{\delta h^J} +\frac{1}{2} \left( \frac{\delta S^1_{IR}}{\delta \phi} \right)^2 - \frac{1}{2} \frac{\delta^2 S^1_{IR}}{\delta \phi^2}.\label{WDW}
\end{split}
\end{equation}

 Let's derive the exact RG equation from the above equation of motion (\ref{WDW}). We use the right hand side and the left hand side of the equality, and use $\delta S^0_{UV} / \delta h + \delta S^0_{IR} / \delta h = 0$, because $S^0$ is classical, purely gravitational part of the action, and the configuration of the metric $h$ should minimize this. (See \cite{VV}) Remember that now the scalar field $\phi$ is quantized and we can't naively consider its classical trajectory. We have the classical gravity background, and we parameterize the physical scale by the overall factor of the metric. By parameterizing $h_{ij}= e^{2a} \delta_{ij}$, "$-a$" can be considered as the parameter for the RG flow.

\begin{equation}
-\frac{d}{da} S^1 = - G^{IJ} \frac{\delta S^0_{UV}}{\delta h^I} \frac{\delta S^1}{\delta h^J}
=-\left( \frac{1}{2} \frac{\delta S^1}{\delta \phi} -\frac{\delta S^1_{UV}}{\delta \phi}  \right) \frac{\delta S^1}{\delta \phi}
 + \left( \frac{1}{2} \frac{\delta^2 S^1}{\delta \phi^2} -\frac{\delta^2 S^1_{UV} }{\delta \phi^2} \right)\label{HERG1}
\end{equation}

%From the first line to the second line, we have used the fact that the radial coordinate $r$ and the scale parameter $a$ can be identified in the pure AdS metric, and that 
Here we have assumed the behavior of the metric $h$ is determined by the leading order of the action $S_{UV}=S^0_{UV} / 16\pi G_N$, together with (\ref{velocity}).

The above expression is a class of the exact RG equation presented in \cite{Morris},\cite{Morris2},\cite{Oliver}. Please remember, in order to identify the above equation as the exact RG, we needed to identify the boundary value of the bulk field as the dynamical field on the boundary CFT as well. At least it is formally true in some cases of AdS/CFT where we can take free boundary condition, there we integrate over all the possible configuration of the boundary value of the bulk field, and it was treated as a dynamical field on the boundary \cite{MukundHRG}, \cite{WittenNeumann}. Moreover, when we put the boundary at the finite radial distance, the boundary condition of the stochastic string \cite{StochasticSon}, \cite{StochasticdeBoer} and graviton \cite{VV} were identified as the dynamical field on the boundary. However, in order to justify it from the field theory, we would need to show that the effective action can be thought of as a classical action, and that the value of the one-point function can be thought of as the value of the field, and that those fields are the fundamental degree of freedom. We will probably need the discussion about the large N field theory as a classical theory \cite{Yaffe}, \cite{Barbon}, and about the formulation of the gauge theory in terms of Wilson loop \cite{largeN}.

%We could simplify the above expression (\ref{HERG1}) by separating the change of the total low energy effective action $S^1$ into the change due to the conformal transformation of the fields, and the change of the functional form of $S^1$ itself due to the conformal transformation. Since the sum of those changes add up to zero, we can write the pure change of the functional form of $S^1$ as,
%\begin{equation}
%-\frac{\partial S^1}{\partial r} = -G^{IJ}\frac{\delta S^0_{UV}}{\delta h^I}\frac{\delta S^1}{\delta h^J}
%- \frac{\delta S^1_{UV}}{\delta \phi}\frac{\delta S^1}{\delta \phi} +\frac{\delta^2 S^1_{UV}}{\delta \phi^2}
%= -\frac{1}{2} \left( \left( \frac{\delta S^1}{\delta \phi} \right)^2 - \frac{\delta^2 S^1}{\delta \phi^2} \right)\label{HERG2}
%\end{equation}
%The first and the second term in the middle expression are the change of $S^1$ itself due to the conformal transformation of $h$ and $\phi$, and the last term comes from the change of the integration measure $\mathcal{D}\phi$. we will explain it in more detail in the next section. The last equality comes from (\ref{HERG1}).

Note that the expression (\ref{HERG1}) can be thought of as a Fokker-Planck equation if we identify $-da$ or $-dr$ as time, and identify probability density as $P= \exp(-S^1)$, and if $S_{UV}$ is nearly constant with respect to $r$ (In the regime of CFT). It implies that it would be also possible to formulate this exact RG equation as a Langevin equation \cite{Gaite}, $ -\partial_a \phi= -\delta S_{UV} / \delta \phi + \xi $, or in terms of the boundary CFT, $S_{UV}$ will be replaced by the generating functional $\Gamma[\phi]$.

%This can be interpreted as the exact RG equation of the CFT, and $S^1$ will be interpreted as the Wilsonian effective action, as the function of the cutoff scale, as we will explain later. The part with $S_{UV}$ can be thought of as the naive scaling from conformal transformation, and together with the change of the metric, it would amount to the change of the coordinate in the CFT. There is one concern, if we could really think boundary value of $\phi$ as the field in the CFT side, at least formally, the integration over $\phi$ would give the partition function from the AdS/CFT though.
%; usually, in AdS/CFT correspondence, we take the value of the bulk field at the boundary to be fixed and treat it as nondynamical source, but here, we rather wish to think of it as a field on the boundary as well. Indeed we need to sum over the boundary condition, which we could interpret as the path integral of the boundary. For that identification, first of all we need to be sure that $\phi$, or the corresponding operator $\mathcal{O}$, is the fundamental degree of freedom, and moreover, we will need some more discussion related to the classical property of the large N limit of the field theory and master field. \cite{Yaffe}, \cite{Barbon}\footnote{It is widely known that in the limit of large N, the quantum fluctuation of the CFT is suppressed with $1/N$. When the quantum effect is taken into account, the field configuration will differ from its classical solution, but the path integral is dominated by one particular configuration, which is called "master field". The master field can be obtained by minimizing an action. (In this sense it can be considered as "classical") The master field action will be identified with the bulk classical supergravity action. If we recover the quantum fluctuation in the CFT around the master field, (which amounts to recovering quantum fluctuation in the bulk,) we expect to get the right partition function.}.As a support, in the setting of \cite{VV}, it is possible to see the fields in the bulk are also dynamical fields in the boundary at finite radial distance. In that picture the above definition of the Wilsonian effective action will be more natural.This simple identification is probably because we didn't make use of open-closed duality, unlike the usual AdS-CFT correspondence. Also in the context of the stochastic string, this sort of straightforward identification between the bulk fundamental string and a particle at the boundary, is made.
%Indeed, in this case, the partition function, where integration over the fields except the boundary surface is made, is identified as the boundary path integral expression. At least, it is true that in the classical limit of the bulk we will attain the correct correspondence, where the boundary condition is one point function $\left< \mathcal{O} \right>$, and classical bulk action is effective action of the CFT. Anyhow, interpreting this bulk field as the boundary field, we obtain one of the class of the exact RG equations which was proposed in \cite{Morris}. It is still necessary to check if this is equivalent to Polchinski RG, following the framework of \cite{Morris3}.  Since this equation has the interpretation as the exact RG, its evolution along $-r$ direction will give us the consistent change of the Wilsonian effective action as we will see. % There is one concern; This is the RG flow equation with respect to the source, while Polchinski type exact RG is equation for the operators. However, it is known that in AdS/CFT correspondence, there is a duality which connects the source and the operator (Strictly speaking, one-point function)\cite{Witten},\cite{Witten2},\cite{Muck}\footnote{This is true for the bulk scalar fields within the mass range, $-d^2/4 < m^2 < -(d^2/4)+1$. } %but we suspect that kind of duality might hold nonperturbatively for general range. Note that we still have a problem with the gravity term, although in some approximation this pert might be neglected.
%, in which we have the duality for the exchange of the source and the operator of the boundary theory, at the same time changing the Energy functional (function of the source, $\Gamma[\phi]$) to the Effective action, $\Gamma[\mathcal{O}]$\footnote{There is a concern that $\mathcal{O}$ is not fundamental degree of freedom, but since its Legendre transform, generating functional is well-defined, effective action also would be proper object to consider.}. Therefore, in the dual perspective, the Generating functional with cutoff will be interpreted as Wisonian effective action, and the RHS will be the derivative with respect to the fields (or, again, strictly speaking, one-point function $\left< \mathcal{O} \right>$). That will be one of the class of the Exact RG equation which was proposed in \cite{Morris}. It is still necessary to check if this is equivalent to Polchinski RG, following the framework of \cite{Morris3}.
%I don't see any problem generalizing this scalar field to general matrix valued fields, since they can be performed exact RG of stochastic quantization, avoiding the gauge invariance problem.


 %I would like to explain why $S^1$, which is sum over the UV part and IR part of the action, could be interpreted as the Wilsonian action. Following \cite{Polchinski}, we can decompose the partition function of the bulk into three parts, UV part, IR part, and on the hypersurface at some finite r. First of all, we can decompose the gravity action into UV part and IR part, since there is no mixing of fields of different $r$ in the gravity action.
 %\begin{equation}
 %\begin{split}
 %\mathcal{Z} &= \int \mathcal{D}\phi_{UV} \mathcal{D} \phi_{IR} \mathcal{D} \phi_r \exp(-S^{Grav}_{UV}) \exp(-S^{Grav}_{IR})\\
 %            &= \int \mathcal{D}\phi_r \Psi_{UV}(r,\phi_r) \Psi_{IR}(r,\phi_r)\\
 %            &= \int \mathcal{D}\phi_r \exp(-S^1(r,\phi_r))\label{WilsonianPartition}
 %\end{split}
 %\end{equation}
%In the second line, we have integrated over IR and UV part of the bulk fields, and in the last line, we have used the definition of the $S^1$, and $\Psi_{UV}$, $\Psi_{IR}$. It is tempting to identify the above expression as the expression of the partition function in the CFT (which is equal to partition function in the AdS) in terms of the Wilsonian effective action, $S^1(r, \phi_r)$. Indeed this action $S^1(r, \phi_r)$ transforms in the way, which is a class of the Exact RG\cite{Morris}\cite{Morris2}\cite{Oliver} with respect to $-r$, as it was obtained from the Wheeler-DeWitt equation (or in the limit of Schrodinger equation) in the bulk. Also, from the AdS/CFT viewpoint, a hypersurface in the AdS is expected to manifest CFT at some physical scale which is specified by $r$, therefore it is natural that Wilsonian effective action "lives" at some r. Also note that, the above definition of the Wilsonian effective action will be consistent with \cite{MukundHRG} in the classical limit of the bulk. However, there is a problem in this identification. In the AdS/CFT correspondence, $S^1$ was rather a Quantum effective action, and $\phi$ was the value of the one point function (Using the duality we explained). Identification of the Wheeler-DeWitt equation and the Exact RG equation is true only if we can interpret $S^1$ as a classical action, and $\phi$ as a classical field. However, there are a line of discussions that, in the limit of Large N, the theory becomes classical in some sense \cite{Yaffe}, \cite{Barbon}, therefore it might enable us to see (\ref{WilsonianPartition}) as the expression of the partition function as the integration over classical fields with the weight of the Wilsonian action.

\section{From the viewpoint of the Exact RG}

Let's see that the expression we have obtained in the last section is actually the exact RG equation in the field theory. First of all, any RG flow equation are obtained from two steps, coarse-graining the microscopic degree of freedom and rescaling. In the continuum field theory, it is known that the coarse-graining step can be simply performed by redefinition of the field\cite{Oliver}. Therefore, we consider the conformal transformation (scale transformation) of the fields in the CFT, and in addition, we perform extra transformation on scalar field as $\mathcal{F}(\phi) dr$, which will correspond to the coarse-graining. The transformation of the fields will be then,
\begin{equation}
\begin{split}
\phi'&= \phi+ \left( \frac{\delta S^1_{UV}}{\delta \phi} + \mathcal{F}(\phi) \right) dr\\
h'   &= h+ \frac{\delta S^0_{UV}}{\delta h}dr
\label{fieldtransf}
\end{split}
\end{equation}

We now assume that the change of the effective action $S^1$ is solely due to the change of the field, and the functional form of $S^1$ doesn't change, as in the regular derivation of the exact RG. We can write down the Wilsonian effective action $S^1(\phi', h')$ with respect to the new fields $\phi'$, $h'$, and by definition it should satisfy
\begin{equation}
\int \mathcal{D}\phi' \exp(-S^1(\phi', h'))=\int \mathcal{D}\phi \exp(-S^1(\phi, h)).
\end{equation}

Expanding the left hand side with respect to $\phi$, $h$, and collecting terms of the order $\mathcal{O}(dr)$, we get
\begin{equation}
- G^{IJ}\frac{\delta S^0_{UV}}{\delta h^I}\frac{\delta S^1}{\delta h^J}-\frac{\delta S^1_{UV}}{\delta \phi}\frac{\delta S^1}{\delta \phi}
-\mathcal{F}\frac{\delta S^1}{\delta \phi} + \frac{\delta^2 S^1_{UV}}{\delta \phi^2}+ \frac{\delta \mathcal{F}}{\delta \phi}=0.
\end{equation}

If we choose $\mathcal{F}= -\frac{1}{2} \frac{ \delta S^1 }{ \delta \phi} $, this expression is precisely what we have obtained from the Wheeler-DeWitt equation in the bulk,(\ref{HERG1}). The form of the field transformation for the coarse graining $\mathcal{F}$ is intuitively reasonable, because, if we see it in the form of the Langevin equation in the previous section, this term corresponds to the noise, therefore it will look like it is washing out microscopic degree of freedom, as any RG should. It was interpreted in \cite{Gaite} that this noise term corresponds to the contribution from irrelevant operators.


%In this section, Let's derive the exact RG equation from the field theory discussion for completeness, following \cite{Oliver}.
%\begin{equation}
%\mathcal{Z}=\int \exp(-S_0 (\phi_0))\mathcal{D}\phi_0,
%\end{equation}
%where $S_0$ is the classical action of the field theory. (Tentatively we wish to identify with the low energy effective action of CFT)\footnote{I am not sure if it is sensible to "quantize" the effective action of the CFT to recover its fluctuation around master field, but the Exact RG equation we got from the gravity is suggestive that we will need to integrate it in order to get the flucuation in CFT (which will correspond to the fluctuation in AdS).} Let's perform the exact RG on this partition function, following (\cite{Oliver},\cite{Morris}). We integrate out the field bit by bit from UV to IR. Let's define "coarse graining function" $b(\phi_0)$ for that purpose, as $b_{dt}(\phi_0)=\phi_0-\Psi(\phi_0) dt$, where $dt= -d\Lambda / \Lambda$, $\Lambda$ being momentum scale of RG. We then define the Wilsonian effective action $S(\phi)$ by
%\begin{equation}
%e^{-S(\phi)}= \int \mathcal{D} \phi_0  \ \delta (\phi-b_{dt}(\phi_0))e^{-S_0(\phi_0)}.
%\end{equation}
%We can easily check it indeed satisfies the necessary condition for the Wilsonian effective action,

%\begin{equation}
%\int \mathcal{D} \phi \ e^{-S(\phi)}= \int \mathcal{D} \phi_0 \ e^{-S_0(\phi_0)}\label{EFGdef}.
%\end{equation}

%Expanding $\mathcal{D}\phi$ and $S(\phi)$ with respect to $\phi_0$, and plugging into (\ref{EFGdef}),

%%\begin{equation}
%%\begin{split}
%%\mathcal{D}\phi &= \det \left( \frac{\mathcal{D} \phi}{\mathcal{D}\phi_0} \right) \mathcal{D}\phi_0=\left(1- \frac{\delta \Psi}{\delta \phi_0} \right)\mathcal{D}\phi_0\\
%%S(\phi) &= S(\phi_0)- \Psi \frac{\delta S_0}{\delta \phi_0}dt\\
%%        &\sim S(\phi_0) - \left( \frac{\delta S_0}{\delta \phi_0} \right) \Psi dt,
%%\end{split}
%%\end{equation}

%we will get the transformation of the Wilsonian effective action as
%\begin{equation}
%S_0(\phi)=S(\phi_0)-\left( \left( \frac{\delta S_0}{\delta \phi_0}\right) \Psi + \frac{\delta \Psi}{\delta \phi_0} \right) dt
%\end{equation}
%That is precisely the equation we have obtained from the gravity side, when we take $\Psi=(1/2)(\delta S^1_0 / \delta \phi_0)- (\delta S^1_{UV} / \delta \phi_0)$. The part with $S_{UV}$ can be understood as naive scaling behavior which is expected from conformal transformation. Therefore, %in addition to the identification that
%the bulk supergravity classical action could be considered as the classical action (for the master field in CFT), with respect to the boundary condition, we can understand that the $S_{UV}$ was the part of the action which has already been integrated out, $S_{IR}$ would be the part before integration, and the sum of them $S$ corresponds to the Wilsonian effective action. Therefore, as we promised, we have shown that the exact RG equation from the gravity side \ref{HERG1} corresponds to the Exact RG of the CFT.

\section{Example: Flow of the double trace operator}

We then still consider the alernative quantization, in which the bulk operators are interpreted as the boundary operator at some RG scale. We mostly follow the discussion in \cite{Polchinski} and \cite{MukundHRG}. We assume that the bulk action is of gaussian form, potential containing only mass term. We take gaussian ansatz for $S_{UV}$ and $S_{IR}$ as a function of the boundary value of the fields $\phi$ and the coordinate $z$. (Note that $z$ can be also understood as the scale parameter from the coordinate dependence of the metric, $h_{ij} \sim z^{-2} \delta_{ij}$.)
The expression of the action for the UV part and the IR part respectively, are
\begin{equation}
\begin{split}
S_{UV}&=C(x,z)+ \frac{1}{2} \int \sqrt{h} \left( F_{UV}(x,z) (\phi-\bar{\phi}(x,z))^2 \right) d^4x \\
S_{IR}&= \frac{1}{2} \int \sqrt{h} \left( F_{IR}(x,z) \phi^2 \right) d^4x .
\end{split}
\end{equation}

We then wish to consider the flow of the double trace operator, which will correspond to the coefficient of the $\phi^2$ term in $S=S_{UV}+S_{IR}$, since, according to the recipe of \cite{dVV}, the sum of the UV and IR action should be identified as the CFT effective action. We will need to see the bulk equation of motion for the evolution of $F=F_{UV}+F_{IR}$, and we can compute $F_{UV}$ and $F_{IR}$ respectively from the Schr\"odinger equation as computed in \cite{Polchinski}. Actually, since we restrict the form of the action to the Gaussian, there is no quantum correction to the double trace operator, (because $(\delta S^2 / \delta \phi^2)$ term will have only constant $\mathcal{O}(\phi^0)$ term) and therefore we can look at Hamilton-Jacobi equation
($\mathcal{H}= z\partial_z S_{UV}=(1/2 \sqrt{h})p^2+(\sqrt{h}/2)( z^2 \delta^{ij} k_i k_j \phi^2 + m^2 \phi^2)$).
%\begin{equation}
%-z \partial_z (S_{UV})=\left( \frac{1}{2\sqrt{h}}\left( \frac{\delta S_{UV}}{\delta \phi} \right)^2
%+\sqrt{h}\left( \frac{1}{2}(\partial_i \phi)^2 -\frac{1}{2} m^2 \phi^2 \right) \right)
%\end{equation}
 Then we will get\footnote{Note that the below result differ from \cite{Polchinski} since their definition of $F$ includes $\sqrt{h}$ as well.}
\begin{equation}
z \partial_z (\frac{1}{2} z^{-d} F_{UV})=\frac{1}{2}(z^{2-d}k^2 + z^{-d}m^2) - \frac{1}{2} z^{-d}F_{UV}^2.
\end{equation}

This equation can be solved by $F_{UV}= z \partial_z \mathcal{O} / \mathcal{O}$, where $\mathcal{O}$ is the Bessel function which solves
\begin{equation}
z \partial_z (z^{-d} z \partial_z \mathcal{O}) = z^{-d} (z^2 k^2 + m^2) \mathcal{O}.
\end{equation}

The above $F_{UV}$, for $kz  \ll  1$, gives
\begin{equation}
F_{UV} \sim \frac{\Delta_- A_{UV}(k)+ \Delta_+ z^{2\nu}}{A_{UV}(k)+ z^{2\nu}}.
\end{equation}
Note that the behavior of $F_{UV}$ for small/large $z$ is as expected, because it should follow the conformal scaling with $\Delta_-$, $\Delta_+$.
On the other hand, we can also obtain $F_{IR}$ from the condition that the solution will not blow up at $z \rightarrow \infty$, and get
\begin{equation}
F_{IR}=- \frac{z \partial_z (z^{d/2}K_{\nu}(kz))}{z^{d/2}K_{\nu}(kz)}.
\end{equation}

When $k \sim 0$, it can be approximated as $F_{IR} \sim -\Delta_-$.
Therefore, if the identification of the total bulk action as the CFT effective action (as \cite{dVV}) is correct, $F$, sum of $F_{UV}$ and $F_{IR}$, should give the right coefficient for the double trace operator, and it is
\begin{equation}
F= F_{IR}+ F_{UV} \sim \frac{(\Delta_+-\Delta_-)z^{2\nu}}{A_{UV}(k)+z^{2\nu}}.
\end{equation}

 Behavior of $F$ at small $z$, $F \sim z^{2\nu}$, is as expected from the conformal dimention of the coupling of the double trace operator. It is also confirmed that this $F \sim F_{UV}-\Delta_-$ follows the same equation as the flow of the coupling of the double trace operator \cite{MukundHRG}. Here we wish to derive the same equation from a different perspective, from the exact RG flow equation that we have derived.
By rewriting (\ref{HERG1}) with respect to $S$ and $S_{IR}$, and concentrating on $\phi^2$ term, we will get an equation

\begin{equation}
z\partial_z F = F^2 -(d-2\Delta_-)F,\label{doubletrace}
\end{equation}
which is the same equation as in \cite{MukundHRG} up to different definition of $F$ by factor of minus sign. We can see that the derivation of the RG flow equation of the double trace operator naturally arises in our case (summing over UV part and IR part of the action). It still remains to be done what will be the modification of the RG flow equation (\ref{doubletrace}) for general $k$.


\section{Discussion}
Therefore, we gave the holographic definition of the Wilsonian effective action, and derived the exact RG flow equation from the Wheeler-DeWitt equation in the bulk, when we can take free boundary condition for the bulk scalar field. The fact, that we could derive the exact RG from the quantum evolution of the bulk, could be convinced in light of stochastic quantization, which also realize the holographic setup; Firstly, It was proposed that the exact RG equation is equivalent to Stochastic RG equation, which looks like Langevin equation, where the scale plays the role of "time". (See \cite{Gaite}\footnote{The apperence of Langevin equation in the context of RG is not limited to Field Theory. See\cite{Oono}, there they also show that Langevin equation appears as the result of RG}. This RG is basically equivalent to Stochastic Quantization.) and also, in the context of the Stochastic quantization, %which we suspect it might be responsible for constructing general holographic models,
the noise of the Langevin equation along the "extra dimension" on the boundary is corresponding to the Quantum fluctuation of the bulk. Summing up those ingredients, we can convince myself that there is a good reason to suspect the quantum evolution of the bulk is related to the "Stochastic RG", or exact RG of the boundary theory. Partial evidence that the Wheeler-DeWitt equation is related to the exact RG equation also comes from \cite{Polchinski}, where they use the Schr\"odinger equation of the bulk to derive the exact RG. Moreover, in the context of the gauge theory of the boundary, the connection between the exact RG and Schwinger-Dyson equation was made \cite{Hirano}, and in addition, it was proposed that the Schwinger-Dyson equation of the boundary is related to the Wheeler-DeWitt equation of the string theory \cite{Gilad}. This gauge theory discussion also suggests the relationship between the Wheeler-DeWitt equation and the exact RG.

This result will hopefully give a better understanding on how exact RG is related to the quantized gravity, and in relating the holographic RG and Polchinski RG. Also, we wish that there will be more understanding of Holography in terms of Stochastic Quantization as exact RG, which was only implied here.
%from the form of the 5d potential
If this identification is established, we will be able to perform the exact RG in a much simpler way, because many miracles, such as no need of gauge-fixing, no need of putting the cutoff, simplification for large N and natural apperance of supersymmetry so on, happens in the stochastic quantization. Note that the same analysis would be applicable to Horava gravity, since they are constructed out of the stochastic quantization, and the stochastic quantization is equivalent to a class of exact RG. It was explained that the structure of stochastic quantization enables the simpler proof of its renormalizability\cite{Susanne}, and we expect that might apply to other general holographic models. Further it was implied that the quantized gravity in the bulk will be written as Langevin equation, which might give more insights into microscopic picture for thermodynamics of spacetime\cite{Jacobson}, \cite{Entropic}.

\section{Update on recent research}
Since this note was originally written more than 10 years ago, there have been a number of recent studies \cite{Sin}\cite{Balasubramanian},\cite{Leigh},\cite{Sathiapalan1},\cite{Sathiapalan2},\cite{Sathiapalan3},\cite{Sathiapalan4},\cite{Sathiapalan5}\cite{Sathiapalan6}, that need to be mentioned.
Some of them \cite{Sin}\cite{Balasubramanian} view the holographic setting in a similar way by splitting to UV part and IR part, although they are based on standard quantization and consider the gravity action after it is integrated over quantum fluctuation.
A series of studies \cite{Sathiapalan1},\cite{Sathiapalan2},\cite{Sathiapalan3},\cite{Sathiapalan4},\cite{Sathiapalan5},\cite{Sathiapalan6} start from the exact RG equation and express transition amplitudes between wave functions of different scale as $\int \mathcal{D}x(t) e^{-S_{bulk}}$, where $x(t)$ denotes fields at different scale $t$, and regard $S_{bulk}$ as the bulk action. They also mention that some of their derivation corresponds to alternative quantization\cite{Sathiapalan2}, and the relationship between Holographic exact RG and Stochastic Quantization\cite{Sathiapalan5}.
The goal of relating the exact RG to the bulk quantum gravity is the same, although they are working on deriving the bulk theory without using explicitly holographic setting.
This note is different in working with alternative quantization and using holographic setting, deriving the exact RG equation on the boundary from the evolution of the wave function in the bulk, dealing with the counterterm.
There are also a number of studies that relate Stochastic Quantization and Holography\cite{SQ1},\cite{SQ2},\cite{SQ3},\cite{SQ4},\cite{SQ5},\cite{SQ6}. In particular, \cite{SQ2},\cite{SQ3},\cite{SQ4},\cite{SQ5},\cite{SQ6} identify radial evolution in the bulk to Stochastic Quantization, which is the exact RG on the boundary, and also derives the evolution of double trace operator, so they are practically very similar to this note. However, their derivation is different from the one in this note, in particular how to subtract the counterterm.
To note, in addition to the research above, we were recently informed that there is another study that mentions the resemblance between the exact RG and Fokker-Planck equation\cite{Hatta}.

\section*{Acknowledgements}
I would like to thank Gilad Lifschitz, Akio Sugamoto and So Katagiri for helpful comments, and those who answered my questions on holography and exact RG.

\begin{thebibliography}{99}

%\cite{de Boer:1999xf}
\bibitem{dVV}
  J.~de Boer, E.~P.~Verlinde and H.~L.~Verlinde,
  ``On the holographic renormalization group,''
  JHEP {\bf 0008} (2000) 003
  [arXiv:hep-th/9912012].
  %%CITATION = JHEPA,0008,003;%%

  %\cite{Verlinde:1999xm}
\bibitem{VV}
  E.~P.~Verlinde and H.~L.~Verlinde,
  ``RG-flow, gravity and the cosmological constant,''
  JHEP {\bf 0005}, 034 (2000)
  [arXiv:hep-th/9912018].
  %%CITATION = JHEPA,0005,034;%%

%\cite{Heemskerk:2010hk}
\bibitem{Polchinski}
  I.~Heemskerk and J.~Polchinski,
  ``Holographic and Wilsonian Renormalization Groups,''
  arXiv:1010.1264 [hep-th].
  %%CITATION = ARXIV:1010.1264;%%

%\cite{Parisi:1980ys}
\bibitem{Parisi}
  G.~Parisi and Y.~s.~Wu,
  ``Perturbation Theory Without Gauge Fixing,''
  Sci.\ Sin.\  {\bf 24}, 483 (1981).
  %%CITATION = SSINA,24,483;%%

%\cite{Damgaard:1988nq}
\bibitem{Damgaard}
  P.H.Damgaard and H.Huffel,
  ``STOCHASTIC QUANTIZATION,''
%\href{http://www.slac.stanford.edu/spires/find/hep/www?irn=1909592}{SPIRES entry}
{\it  SINGAPORE, SINGAPORE: WORLD SCIENTIFIC (1988) 496p},

%\cite{Fukuma:2002sb}
\bibitem{Sakai}
  M.~Fukuma, S.~Matsuura and T.~Sakai,
  ``Holographic renormalization group,''
  Prog.\ Theor.\ Phys.\  {\bf 109}, 489 (2003)
  [arXiv:hep-th/0212314].
  %%CITATION = PTPKA,109,489;%%

%\cite{Morris:1993qb}
\bibitem{Morris}
  T.~R.~Morris,
  ``The Exact renormalization group and approximate solutions,''
  Int.\ J.\ Mod.\ Phys.\  A {\bf 9}, 2411 (1994)
  [arXiv:hep-ph/9308265].
  %%CITATION = IMPAE,A9,2411;%%,

%\cite{Morris:1993qb}
\bibitem{Morris2}
    T.~R.~Morris,
  ``A manifestly gauge invariant exact renormalization group,''
  arXiv:hep-th/9810104.
  %%CITATION = HEP-TH/9810104;%%,

%\cite{Morris:1993qb}
\bibitem{Morris3}
    J.~I.~Latorre and T.~R.~Morris,
  ``Exact scheme independence,''
  JHEP {\bf 0011}, 004 (2000)
  [arXiv:hep-th/0008123].
  %%CITATION = JHEPA,0011,004;%%

%\cite{Rosten:2010vm}
\bibitem{Oliver}
  O.~J.~Rosten,
  ``Fundamentals of the Exact Renormalization Group,''
  arXiv:1003.1366 [hep-th].
  %%CITATION = ARXIV:1003.1366;%%

  %\cite{Klebanov:1999tb}
\bibitem{Witten2}
  I.~R.~Klebanov and E.~Witten,
  ``AdS/CFT correspondence and symmetry breaking,''
  Nucl.\ Phys.\  B {\bf 556}, 89 (1999)
  [arXiv:hep-th/9905104].
  %%CITATION = NUPHA,B556,89;%%

%\cite{Mueck:1999kk}
\bibitem{Muck}
  W.~Mueck and K.~S.~Viswanathan,
  ``Regular and irregular boundary conditions in the AdS/CFT  correspondence,''
  Phys.\ Rev.\  D {\bf 60}, 081901 (1999)
  [arXiv:hep-th/9906155].
  %%CITATION = PHRVA,D60,081901;%%

%\cite{Witten:2001ua}
\bibitem{Witten}
  E.~Witten,
  ``Multi-trace operators, boundary conditions, and AdS/CFT correspondence,''
  arXiv:hep-th/0112258.
  %%CITATION = HEP-TH/0112258;%%

%\cite{Witten:2003ya}
\bibitem{WittenNeumann}
  E.~Witten,
  ``SL(2,Z) action on three-dimensional conformal field theories with Abelian symmetry,''
  arXiv:hep-th/0307041.
  %%CITATION = HEP-TH/0307041;%%

%\cite{Faulkner:2010jy}
\bibitem{MukundHRG}
  T.~Faulkner, H.~Liu, M.~Rangamani,
  ``Integrating out geometry: Holographic Wilsonian RG and the membrane paradigm,''
  [arXiv:1010.4036 [hep-th]].

%\cite{Aharony:2010ay}
\bibitem{MukundOfer}
  O.~Aharony, D.~Marolf and M.~Rangamani,
  ``Conformal field theories in anti-de Sitter space,''
  JHEP {\bf 1102}, 041 (2011)
  [arXiv:1011.6144 [hep-th]].
  %%CITATION = JHEPA,1102,041;%%

%\cite{Marolf:2006nd}
\bibitem{Ross}
  D.~Marolf and S.~F.~Ross,
  ``Boundary conditions and new dualities: Vector fields in AdS/CFT,''
  JHEP {\bf 0611}, 085 (2006)
  [arXiv:hep-th/0606113].
  %%CITATION = JHEPA,0611,085;%%

%\cite{deBoer:2008gu}
\bibitem{StochasticdeBoer}
  J.~de Boer, V.~E.~Hubeny, M.~Rangamani {\it et al.},
  ``Brownian motion in AdS/CFT,''
  JHEP {\bf 0907}, 094 (2009).
  [arXiv:0812.5112 [hep-th]].
  
\bibitem{StochasticSon}
  D.~T.~Son, D.~Teaney,
  ``Thermal Noise and Stochastic Strings in AdS/CFT,''
  JHEP {\bf 0907}, 021 (2009).
  [arXiv:0901.2338 [hep-th]].

%\cite{Polchinski:1983gv}
\bibitem{PolchinskiRG}
  J.~Polchinski,
  ``Renormalization And Effective Lagrangians,''
  Nucl.\ Phys.\  B {\bf 231}, 269 (1984).
  %%CITATION = NUPHA,B231,269;%%

%\cite{Berenstein:2002sa}

%\cite{Lifschytz:1994pd}
\bibitem{GiladSchwinger}
  G.~Lifschytz, S.~D.~Mathur and M.~Ortiz,
  ``A Note on the semiclassical approximation in quantum gravity,''
  Phys.\ Rev.\  D {\bf 53}, 766 (1996)
  [arXiv:gr-qc/9412040].
  %%CITATION = PHRVA,D53,766;%%

%\cite{Nishioka:2009iq}
\bibitem{Nishioka}
  T.~Nishioka,
  ``Horava-Lifshitz Holography,''
  Class.\ Quant.\ Grav.\  {\bf 26}, 242001 (2009)
  [arXiv:0905.0473 [hep-th]].
  %%CITATION = CQGRD,26,242001;%%

%\cite{Orlando:2009en}
\bibitem{Susanne}
  D.~Orlando and S.~Reffert,
  ``On the Renormalizability of Horava-Lifshitz-type Gravities,''
  Class.\ Quant.\ Grav.\  {\bf 26}, 155021 (2009)
  [arXiv:0905.0301 [hep-th]].
  %%CITATION = CQGRD,26,155021;%%

%\cite{Yaffe:1981vf}
\bibitem{Yaffe}
  L.~G.~Yaffe,
  ``Large N Limits As Classical Mechanics,''
  Rev.\ Mod.\ Phys.\  {\bf 54}, 407 (1982).
  %%CITATION = RMPHA,54,407;%%

%\cite{Barbon:2002xk}
\bibitem{Barbon}
  J.~L.~F.~Barbon,
  ``Multitrace AdS/CFT and master field dynamics,''
  Phys.\ Lett.\  B {\bf 543}, 283 (2002)
  [arXiv:hep-th/0206207].
  %%CITATION = PHLTA,B543,283;%%

%\cite{Makeenko:1999hq}
\bibitem{largeN}
  Y.~Makeenko,
  ``Large-N gauge theories,''
  arXiv:hep-th/0001047.
  %%CITATION = HEP-TH/0001047;%%

%\cite{Gaite}
\bibitem{Gaite}
  J.~C.~Gaite,
  ``Stochastic formulation of the renormalization group: Supersymmetric
  structure and topology of the space of couplings,''
  J.\ Phys.\ A  {\bf 37}, 10409 (2004)
  [arXiv:hep-th/0404212].
  %%CITATION = JPAGB,A37,10409;%%

\bibitem{Oono}
Some implications of renormalization group theoretical ideas to statistics
Physica D: Nonlinear Phenomena, Volume 205, Issue 1-4, June 2005, Pages 207-214
Rajaram, S.; Taguchi, Y.h.; Oono, Y.

%\cite{Hirano:1999ay}
\bibitem{Hirano}
  S.~Hirano,
  ``Exact renormalization group and loop equation,''
  Phys.\ Rev.\  D {\bf 61}, 125011 (2000)
  [arXiv:hep-th/9910256].
  %%CITATION = PHRVA,D61,125011;%%

  %\cite{Lifschytz:2000bj}
\bibitem{Gilad}
  G.~Lifschytz and V.~Periwal,
  ``Schwinger-Dyson = Wheeler-DeWitt: Gauge theory observables as bulk
  operators,''
  JHEP {\bf 0004}, 026 (2000)
  [arXiv:hep-th/0003179].
  %%CITATION = JHEPA,0004,026;%%

%\cite{Jacobson:1995ab}
\bibitem{Jacobson}
  T.~Jacobson,
  ``Thermodynamics of space-time: The Einstein equation of state,''
  Phys.\ Rev.\ Lett.\  {\bf 75}, 1260 (1995)
  [arXiv:gr-qc/9504004].
  %%CITATION = PRLTA,75,1260;%%

%\cite{Verlinde:2010hp}
\bibitem{Entropic}
  E.~P.~Verlinde,
  ``On the Origin of Gravity and the Laws of Newton,''
  arXiv:1001.0785 [hep-th].
  %%CITATION = ARXIV:1001.0785;%%

%\cite{Li:2000ec}
%\bibitem{Miao}
%  M.~Li,
%  ``A note on relation between holographic RG equation and Polchinski's RG
%  equation,''
%  Nucl.\ Phys.\  B {\bf 579}, 525 (2000)
%  [arXiv:hep-th/0001193].
  %%CITATION = NUPHA,B579,525;%%

%\cite{Sin:2011yh}
\bibitem{Sin}
S.~J.~Sin and Y.~Zhou,
``Holographic Wilsonian RG Flow and Sliding Membrane Paradigm,''
JHEP \textbf{05}, 030 (2011)
doi:10.1007/JHEP05(2011)030
[arXiv:1102.4477 [hep-th]].

%\cite{Balasubramanian:2012hb}
\bibitem{Balasubramanian}
V.~Balasubramanian, M.~Guica and A.~Lawrence,
``Holographic Interpretations of the Renormalization Group,''
JHEP \textbf{01}, 115 (2013)
doi:10.1007/JHEP01(2013)115
[arXiv:1211.1729 [hep-th]].

%\cite{Leigh:2014qca}
\bibitem{Leigh}
R.~G.~Leigh, O.~Parrikar and A.~B.~Weiss,
``Exact renormalization group and higher-spin holography,''
Phys. Rev. D \textbf{91}, no.2, 026002 (2015)
doi:10.1103/PhysRevD.91.026002
[arXiv:1407.4574 [hep-th]].

%\cite{Sathiapalan:2017frk}
\bibitem{Sathiapalan1}
B.~Sathiapalan and H.~Sonoda,
``A Holographic form for Wilson's RG,''
Nucl. Phys. B \textbf{924}, 603-642 (2017)
doi:10.1016/j.nuclphysb.2017.09.018
[arXiv:1706.03371 [hep-th]].

%\cite{Sathiapalan:2019zex}
\bibitem{Sathiapalan2}
B.~Sathiapalan and H.~Sonoda,
``Holographic Wilson's RG,''
Nucl. Phys. B \textbf{948}, 114767 (2019)
doi:10.1016/j.nuclphysb.2019.114767
[arXiv:1902.02486 [hep-th]].

%\cite{Sathiapalan:2020cff}
\bibitem{Sathiapalan3}
B.~Sathiapalan,
``Holographic RG and Exact RG in O(N) Model,''
Nucl. Phys. B \textbf{959}, 115142 (2020)
doi:10.1016/j.nuclphysb.2020.115142
[arXiv:2005.10412 [hep-th]].

%\cite{Dharanipragada:2020ibl}
\bibitem{Sathiapalan4}
P.~Dharanipragada, S.~Dutta and B.~Sathiapalan,
``Bulk gauge fields and holographic RG from exact RG,''
JHEP \textbf{23}, 174 (2020)
doi:10.1007/JHEP02(2023)174
[arXiv:2201.06240 [hep-th]].

%\cite{Dharanipragada:2022yxw}
\bibitem{Sathiapalan5}
P.~Dharanipragada, S.~Dutta and B.~Sathiapalan,
``Aspects of the map from exact RG to holographic RG in AdS and dS,''
Mod. Phys. Lett. A \textbf{37}, no.37n38, 2250235 (2022)
doi:10.1142/S0217732322502352
[arXiv:2301.13605 [hep-th]].

\bibitem{Sathiapalan6}
P.~Dharanipragada and B.~Sathiapalan,
``Holographic RG from ERG: Locality and General Coordinate Invariance in the Bulk,''
[arXiv:2306.07442 [hep-th]].

\bibitem{SQ1}
D.~S.~Mansi, A.~Mauri and A.~C.~Petkou,
``Stochastic Quantization and AdS/CFT,''
Phys. Lett. B \textbf{685}, 215-221 (2010)
doi:10.1016/j.physletb.2010.01.033
[arXiv:0912.2105 [hep-th]].

\bibitem{SQ2}
J.~H.~Oh and D.~P.~Jatkar,
``Stochastic quantization and holographic Wilsonian renormalization group,''
JHEP \textbf{11}, 144 (2012)
doi:10.1007/JHEP11(2012)144
[arXiv:1209.2242 [hep-th]].

\bibitem{SQ3}
D.~P.~Jatkar and J.~H.~Oh,
``Stochastic quantization of conformally coupled scalar in AdS,''
JHEP \textbf{10}, 170 (2013)
doi:10.1007/JHEP10(2013)170
[arXiv:1305.2008 [hep-th]].

%\cite{Oh:2021bxx}
\bibitem{SQ4}
J.~H.~Oh,
``First-order formalism of holographic Wilsonian renormalization group: Langevin equation,''
J. Korean Phys. Soc. \textbf{79}, no.10, 903-917 (2021)
doi:10.1007/s40042-021-00320-x
[arXiv:2110.05013 [hep-th]].

\bibitem{SQ5}
G.~Kim, J.~s.~Chae, W.~Shin and J.~H.~Oh,
``Stochastic quantization and holographic Wilsonian renormalization group of scalar theory with generic mass, self-interaction and multiple trace deformation,''
Int. J. Mod. Phys. A \textbf{38}, no.21, 2350114 (2023)
doi:10.1142/S0217751X23501142
[arXiv:2305.18920 [hep-th]].

\bibitem{SQ6}
J.~H.~Lee and J.~H.~Oh,
``Stochastic quantization and holographic Wilsonian renormalization group of conformally coupled scalar in AdS$_{4}$,''
J. Korean Phys. Soc. \textbf{83}, no.9, 665-674 (2023)
doi:10.1007/s40042-023-00926-3
[arXiv:2308.10010 [hep-th]].

%\cite{Hatta:2001ui}
\bibitem{Hatta}
Y.~Hatta and T.~Kunihiro,
``Renormalization group method applied to kinetic equations: Roles of initial values and time,''
Annals Phys. \textbf{298}, 24-57 (2002)
doi:10.1006/aphy.2002.6234
[arXiv:hep-th/0108159 [hep-th]].


\end{thebibliography}

\end{document}
question: effective action of the bulk and the boundary should be the same, how come we can normally compute the CFT effective action from classical gravity (in the CFT sense. What are we ignoring in the CFT when we ignore quantum fluctuation in the bulk? Maybe it is related to the fact Effective action is a quantity which is related to RG flow, and we are ignoring stochasticity in the RG flow equation and it affects the effective action?)
